\subsection{Lexical Closures}

A closure is an inner function whose code references free variables from an enclosing scope, preserved in heap-allocated cell objects after the outer frame unwinds. The outer frame is not kept alive as an active call record; only the referenced bindings survive via \texttt{\_\_closure\_\_}.

\subsubsection{Closures, Lexical Enclosures, and Free Variable Tracking Invariants}

When the compiler detects that an inner function reads an outer local, it promotes that name to a \texttt{cell}: a \texttt{PyCellObject} holding a \texttt{PyObject*} pointer. The outer code object's \texttt{co\_cellvars} lists names stored in cells; the inner code object's \texttt{co\_freevars} lists names loaded through cells. At runtime, \texttt{LOAD\_DEREF} reads \texttt{cell\_contents}; \texttt{STORE\_DEREF} writes it.

%<<<PROGRAM:programs/lex01>>>


Each call to a factory creates separate function objects with separate cell tuples holding distinct object references. Mutable cells share one object across invocations of the same closure (\texttt{items.append(...)}); immutable rebinding requires \texttt{nonlocal} to update the cell pointer.

Reference counting on cells keeps enclosed objects alive until the closure itself is collected. This is explicit heap indirection---not a copy of the outer frame, not a global variable.

Pair \texttt{co\_freevars} with \texttt{\_\_closure\_\_} by position to inspect preserved state. A function with no free variables has \texttt{\_\_closure\_\_} is \texttt{None}.

The lifespan shift: \texttt{outer()} returns, its frame is deallocated, but \texttt{fn()} still reads \texttt{message} through the cell whose reference count remains nonzero because \texttt{fn} holds it in \texttt{\_\_closure\_\_}.
